\documentclass[10pt,letterpaper]{article}
\usepackage{cornell-notes}

\title{Clause 22.14.08.03: Class Template Basic Format Args}
\author{}
\date{}
\setDocTitle{Clause 22.14.08.03: Class Template Basic Format Args}
\setDocSubtitle{C++ 2024}
\setDocOwner{C++ 2024 Cornell Notes}
\setDocDate{\today}
\setCornellCollection{C++ 2024 Cornell Notes}
\setCornellUnitType{Clause}
\setCornellUnitNumber{22.14.08.03}
\setCornellUnitTitle{Class Template Basic Format Args}
\hypersetup{pdftitle={Clause 22.14.08.03: Class Template Basic Format Args},pdfsubject={Cornell notes},pdfauthor={}}

\begin{document}
\maketitle
\begin{CornellOverviewBox}{Overview}
\textbf{Classification:} Standard library - Normative\\[2pt]
\textbf{Learning objective:} Understand the rules, terminology, and semantic consequences associated with Class template basic\_format\_args.
\end{CornellOverviewBox}\begin{CornellSummaryBox}{Summary}
{\sffamily\bfseries\color{Accent} Summary}\par\vspace{3pt}Section 22.14.8.3 focuses on Class template basic\_format\_args. Apply the specified interface together with its mandates, effects, result conditions, exception behavior, and complexity constraints. When applying it, preserve the distinction between mandatory requirements, recommendations, permissions, and behavior deliberately left undefined, unspecified, or implementation-defined.
\end{CornellSummaryBox}

\end{document}
